\documentclass{beamer}

\usepackage[utf8]{inputenc}
\usepackage[T1]{fontenc}
\usepackage{microtype}
\usepackage[german]{babel}
\usepackage{hyperref}
\usepackage{lmodern}
\usepackage{wasysym}
\usepackage{textpos}
\usepackage{ulem}
\usepackage{epstopdf}
\usepackage[pageofpages=von,
            titlepagelogo=resource/logo]{beamerthemeUhh}
\usepackage{beamercolorthemeuhh}

\usepackage{tikz}
\usepackage{setspace}

\newenvironment{changemargin}[2]{%
\begin{list}{}{%
\setlength{\topsep}{0pt}%
\setlength{\leftmargin}{#1}%
\setlength{\rightmargin}{#2}%
\setlength{\listparindent}{\parindent}%
\setlength{\itemindent}{\parindent}%
\setlength{\parsep}{\parskip}%
}%
\item[]}{\end{list}}

\logo{\includegraphics[height=0.125\paperheight]{resource/logo-fbi}}

\setbeamertemplate{frametitle continuation}[from second][]
\usefonttheme{professionalfonts}

%%%%%
\newcommand{\myName}{Felix Favre}
\newcommand{\myEmail}{1favre@informatik.uni-hamburg.de}
\newcommand{\myTitle}{Validierung von dynamisch geladenem Javascript Code in Webapplikationen}
\newcommand{\myKeywords}{}
\institute{Universität Hamburg -- Fachbereich Informatik -- Abschlussarbeiten Seminar}
\date{\today}


\title{\myTitle}
\author{\myName}
\keywords{\myKeywords}
\subject{}

\definecolor{uhh}{RGB}{226,0,26}

\hypersetup{
	pdftitle={\myTitle},
	pdfauthor={\myName},
	pdfkeywords={\myKeywords},
	bookmarksnumbered=true,
	bookmarksopen=true,
	bookmarksopenlevel=1,
	colorlinks=true,
	linkbordercolor={1 1 1},
	urlcolor=uhh,
	anchorcolor=black,
	linkcolor=black,
	citecolor=black,
	filecolor=black,
	menucolor=black
}

\pdfcompresslevel=9
\pdfimageresolution=72
\pdfpkresolution=72

\newcommand{\maxFrameImage}[1]{
\begingroup
\setbeamercolor{background canvas}{bg=black}
\begin{frame}[plain]
\begin{changemargin}{-1cm}{-1cm}
\begin{center}
\includegraphics[width=\paperwidth,height=\paperheight,keepaspectratio]
{#1}
\end{center}
\end{changemargin}
\end{frame}
\endgroup
}

\begin{document}

\pdfbookmark[2]{Titelseite}{title}
{
    \usebackgroundtemplate{\includegraphics[width=\paperwidth]{resource/uhh-bg.png}}
    \frame{
        \titlepage
    }
}

\begin{frame}{Was geplant war}
  \begin{itemize}
    \item Git Repository aufsetzen (GitHub)
    \item Redmine aufsetzen
    \item Deckblatt designen
    \item Gliederung schreiben
    \item Literatur raussuchen
  \end{itemize}
\end{frame}

\begin{frame}{Was geplant war}
  \begin{itemize}
    \item \sout{Git Repository aufsetzen (GitHub)}
    \item Redmine aufsetzen
    \item Deckblatt designen
    \item Gliederung schreiben
    \item Literatur raussuchen
  \end{itemize}
\end{frame}

\begin{frame}{Was geplant war}
  \begin{itemize}
    \item \sout{Git Repository aufsetzen (GitHub)}
    \item \sout{Redmine aufsetzen}
    \item Deckblatt designen
    \item Gliederung schreiben
    \item Literatur raussuchen
  \end{itemize}
\end{frame}

\begin{frame}{Was geplant war}
  \begin{itemize}
    \item \sout{Git Repository aufsetzen (GitHub)}
    \item \sout{Redmine aufsetzen}
    \item \sout{Deckblatt designen}
    \item Gliederung schreiben
    \item Literatur raussuchen
  \end{itemize}
\end{frame}

\begin{frame}{Was geplant war}
  \begin{itemize}
    \item \sout{Git Repository aufsetzen (GitHub)}
    \item \sout{Redmine aufsetzen}
    \item \sout{Deckblatt designen}
    \item \sout{Gliederung schreiben}
    \item Literatur raussuchen
  \end{itemize}
\end{frame}

\begin{frame}{Was geplant ist}
  \begin{itemize}
    \item Motivation ausformulieren
    \item Gliederung weiter verfeinern
    \item Literatur raussuchen
    \item Mehr schreiben
  \end{itemize}
\end{frame}

\end{document}
